\documentclass[11pt]{article}

\usepackage[a4paper,margin=1in]{geometry}
\usepackage{amsmath,amssymb,amsthm,mathtools}
\usepackage[hidelinks]{hyperref}
\usepackage{microtype}

\newtheorem{theorem}{Theorem}[section]
\newtheorem{lemma}[theorem]{Lemma}
\newtheorem{proposition}[theorem]{Proposition}
\newtheorem{corollary}[theorem]{Corollary}
\theoremstyle{remark}
\newtheorem{remark}[theorem]{Remark}

\newcommand{\R}{\mathbb{R}}
\newcommand{\K}{\mathbb{K}}
\newcommand{\OO}{\mathcal{O}}
\newcommand{\length}{\operatorname{length}}
\newcommand{\val}{\operatorname{val}}
\newcommand{\supp}{\operatorname{supp}}
\newcommand{\one}{\mathbf{1}}
\newcommand{\unionword}[2]{#1#2}

\title{Uniform Ingleton Bounds for Nonnegative Lorentzian Quadratics}
\author{guangxiangdebizi}
\date{27 August 2026}

\begin{document}
\maketitle

\begin{abstract}
For a nonzero homogeneous quadratic $f$ with nonnegative coefficients and a
subset $S$ of its variables, put
\[
 V_\varepsilon(S)=f(\one_S+\varepsilon\one_E),\qquad \varepsilon>0.
\]
We prove that if $f$ is Lorentzian, then the multiplicative Ingleton ratio
formed from the ten values $V_\varepsilon(S)$ is bounded below by a universal
positive constant.  The constant is independent of the number of variables,
the polynomial, four arbitrary and possibly overlapping subsets, and the
regularization parameter.  The proof reduces to sixteen membership atoms,
uses semialgebraic curve selection to tropicalize a hypothetical vanishing
sequence, realizes the resulting rank-two valuated matroid over a discretely
valued field, and places all ten tropical evaluations in a single
finite-length module.  The ordinary Ingleton inequality for module lengths
then rules out positive ratio valuation.  We also construct an exact
five-vector determinantal family showing that the optimal constant is at most
\[
 \frac{-107+51\sqrt{17}}{128}=0.806862397707\ldots.
\]
The proof is non-effective and does not determine the optimal constant.
\end{abstract}

\section{Statement}

Let $E$ be a finite set and let
\[
 f(x)=\sum_{|\alpha|=2}c_\alpha x^\alpha,
 \qquad c_\alpha\geq 0,
\]
be nonzero.  In degree two we call $f$ \emph{Lorentzian} when its Hessian has
at most one positive eigenvalue.  This agrees with the degree-two part of the
standard definition of Lorentzian polynomials
\cite{BrandenHuh2020}.

For $S\subseteq E$ and $\varepsilon>0$, define
\begin{equation}
 V_{f,\varepsilon}(S)=f(\one_S+\varepsilon\one_E).
 \label{eq:evaluation}
\end{equation}
Every coordinate in \eqref{eq:evaluation} is positive, so every such value is
positive.  Given four arbitrary subsets $A,B,C,D\subseteq E$, write $AB$ for
$A\cup B$, and similarly for longer strings.  Define
\begin{equation}
 \mathcal R_f(\varepsilon;A,B,C,D)=
 \frac{V(AB)V(AC)V(AD)V(BC)V(BD)}
 {V(A)V(B)V(ABC)V(ABD)V(CD)}.
 \label{eq:ratio}
\end{equation}
where $V$ abbreviates $V_{f,\varepsilon}$.  Let
\begin{equation}
 c_2^\star=
 \inf_{E,f,A,B,C,D,\varepsilon}
 \mathcal R_f(\varepsilon;A,B,C,D),
 \label{eq:optimal-constant}
\end{equation}
where the infimum ranges over exactly the data just specified: finite $E$,
nonzero nonnegative Lorentzian quadratics $f$, arbitrary subsets
$A,B,C,D\subseteq E$, and $\varepsilon>0$.

\begin{theorem}[Uniform degree-two bound]
\label{thm:main}
The optimal constant in \eqref{eq:optimal-constant} satisfies $c_2^\star>0$.
\end{theorem}

The proof gives existence but no effective value.  The following construction
gives a strict upper bound.

\begin{proposition}[Exact upper-bound family]
\label{prop:upper}
The best possible constant in Theorem~\ref{thm:main} satisfies
\[
 0<c_2^\star\leq \frac{-107+51\sqrt{17}}{128}.
\]
\end{proposition}

Classical Ingleton says that a rank function obtained from linear subspaces
satisfies an additive five-against-five inequality
\cite{Ingleton1971}.  Theorem~\ref{thm:main} is not a direct application:
the coefficients of $f$ and the regularizer may approach zero at unrelated
rates.  Controlling this simultaneous degeneration is the main point.

\section{Fixed-dimensional reduction}

Let $\mathcal L=\{\mathsf A,\mathsf B,\mathsf C,\mathsf D\}$ be a set of
four labels.  Membership in the four ground-set subsets gives each variable a
pattern $\omega\subseteq\mathcal L$: for example,
$\mathsf A\in\omega$ means that the variable belongs to the ground-set subset
$A$.  Replace every variable with pattern $\omega$ by one variable
$y_\omega$.  There are at most $2^4=16$ such variables, including the outside
pattern $\varnothing$.  Denote the resulting quadratic by $q$.
We abbreviate label subsets by strings, for example
$\mathsf{AB}=\{\mathsf A,\mathsf B\}$.

This substitution preserves all ten evaluations in \eqref{eq:ratio}.  If its
matrix is denoted by $L$, then
\[
 \nabla^2q=L^{\mathsf T}(\nabla^2f)L.
\]
Therefore $q$ is Lorentzian.  We may henceforth work in sixteen fixed
variables, allowing some variables to be inactive.

The ratio is unchanged by multiplying $q$ by a positive scalar, so normalize
\begin{equation}
 q(\one)=\sum_\alpha c_\alpha=1.
 \label{eq:normalization}
\end{equation}
The nonnegative coefficient simplex satisfying \eqref{eq:normalization},
intersected with the closed condition that the Hessian has at most one
positive eigenvalue, is compact and semialgebraic.

Put
\[
 \delta=\frac{\varepsilon}{1+\varepsilon}\in(0,1),\qquad
 u_J(\omega)=
 \begin{cases}
 1,&\omega\cap J\neq\varnothing,\\
 \delta,&\omega\cap J=\varnothing,
 \end{cases}
\]
for a label set $J\subseteq\mathcal L$.  Let
$U_J=\bigcup_{\mathsf X\in J}X$ be the corresponding union of ground-set
subsets.  Thus $\omega\cap J$ is an intersection of two label sets, not an
intersection with the ground set.  Homogeneity gives
\[
 V_{q,\varepsilon}(U_J)=(1+\varepsilon)^2q(u_J),
\]
and the common factor cancels in \eqref{eq:ratio}.  By
\eqref{eq:normalization},
\begin{equation}
 \delta^2\leq q(u_J)\leq 1.
 \label{eq:delta-bounds}
\end{equation}
Consequently a sequence on which the ratio tends to zero must have
$\delta\to0$.

\section{A valuated rank-two limit}

Assume, for contradiction, that the infimum in
Theorem~\ref{thm:main} is zero.  The graph of the ratio over the compact
normalized coefficient space and $0<\delta\leq1$ is semialgebraic.
The Nash curve selection lemma
\cite[Proposition~8.1.13]{BochnakCosteRoy1998} gives a Nash arc approaching a
boundary point with ratio zero.  Each nonzero coordinate of a Nash arc is
algebraic over $\R(t)$; applying the classical Newton--Puiseux theorem to the
finitely many coordinate functions and taking one common finite ramification
therefore gives an arc indexed by $t\downarrow0$ such
that
\begin{align*}
 c_\alpha(t)&=a_\alpha t^{p_\alpha}+\text{higher terms},
 &a_\alpha&>0,\\
 \delta(t)&=a_\delta t^\lambda+\text{higher terms},
 &a_\delta&>0,\quad\lambda>0,
\end{align*}
for every coefficient not identically zero.  Identically zero coefficients
are assigned valuation $+\infty$.  The Hessian signature condition is
semialgebraic, so the arc defines a Lorentzian polynomial over the ordered
field of real convergent Puiseux series.  Since the ratio tends to zero along
the arc,
\begin{equation}
 \val_t\mathcal R_q>0.
 \label{eq:positive-valuation}
\end{equation}

For every atom $i$, introduce two clones $i^0,i^1$, and polarize
\[
 q(y)=\sum_i c_{ii}y_i^2+\sum_{i<j}c_{ij}y_iy_j
\]
to
\begin{equation}
 \widetilde q(z)=
 \sum_i c_{ii}z_{i^0}z_{i^1}
 +\sum_{i<j}\frac{c_{ij}}4
 (z_{i^0}+z_{i^1})(z_{j^0}+z_{j^1}).
 \label{eq:polarization}
\end{equation}
Polarization preserves the Lorentzian property
\cite{BrandenHuh2020}.  Let $p_I$ be the valuation of the coefficient of the
clone pair $I$.  Br\"and\'en and Huh prove that coefficient valuations of
Lorentzian polynomials are $M$-convex
\cite[Corollary~3.28]{BrandenHuh2020}.  In the multiaffine degree-two case this
is exactly the valuated-basis-exchange condition for a rank-two valuated
matroid.  Equivalently, for every four distinct clones $i,j,k,l$,
\begin{equation}
 \min\{p_{ij}+p_{kl},p_{ik}+p_{jl},p_{il}+p_{jk}\}
 \quad\text{is attained at least twice}.
 \label{eq:tropical-plucker}
\end{equation}
Infinite entries are allowed, and the finite pairs form the bases of a
rank-two matroid.

For a label set $J\subseteq\mathcal L$, let $\widetilde E_J$ contain both
clones of every atom $\omega$ satisfying $\omega\cap J\neq\varnothing$, and
define
\begin{equation}
 r_\lambda(J)=
 \max_{I}\{\lambda|I\cap\widetilde E_J|-p_I\},
 \label{eq:regularized-rank}
\end{equation}
where the maximum ranges over finite clone pairs.  All leading coefficients
are nonnegative, so no cancellation occurs.  Equations
\eqref{eq:polarization} and \eqref{eq:regularized-rank} give
\begin{equation}
 \val_t q(u_J(t))=2\lambda-r_\lambda(J).
 \label{eq:evaluation-valuation}
\end{equation}

\section{Rank two over a discrete valuation ring}

We use the special fact that every rank-two valuated matroid is realizable.

\begin{lemma}[Rank-two realization]
\label{lem:realization}
After multiplying all finite $p_I$ by one positive integer, there are a
discretely valued field $\K$ with infinite residue field, valuation ring
$\OO$, uniformizer $\pi$, and columns $v_e\in\K^2$ such that
\[
 \val\det(v_e,v_{e'})=p_{ee'}.
\]
An infinite value is represented by an identically zero determinant.
\end{lemma}

\begin{proof}
For finite uniform support, \eqref{eq:tropical-plucker} is the four-point
condition for a rational weighted tree after adding leaf potentials.  Root
the tree.  At each branching choose distinct residue constants and construct
Puiseux coordinates $x_e$ along the root-to-leaf paths.  The first branch at
which two leaves differ gives
\[
 \val(x_e-x_{e'})=h(e,e'),
\]
where $h(e,e')$ is their last-common-ancestor depth.  Suitable leaf scalars
then make $v_e=\pi^{b_e}(1,x_e)$ realize $p_{ee'}$.  This is the usual
rank-two tree model for the tropical Grassmannian
\cite{SpeyerSturmfels2004}.

For nonuniform support, first remove loops and pass to the parallel classes
of the underlying rank-two matroid.  Choose a representative $e_P$ of every
class $P$.  If $e\in P$ and $f$ is outside $P$, the difference
\[
 p_{ef}-p_{e_Pf}
\]
is independent of the choice of $f$: this is the equality of the two finite
sums in \eqref{eq:tropical-plucker} applied to two elements of $P$ and two
outside elements.  Denote the difference by $s_e$.  After treating the
vacuous two- and three-element cases directly, one obtains for distinct
classes $P,Q$
\[
 p_{ef}=\bar p_{PQ}+s_e+s_f
 \qquad(e\in P, f\in Q).
\]
Realize the simple quotient valuation $\bar p$ by the tree construction and
replace the representative column of $P$ by its multiples $\pi^{s_e}v_{e_P}$.
Elements in one parallel class are then exact scalar multiples, so their
determinants are zero and have value $+\infty$.  Loops are represented by
zero columns.  A finite ramification clears the rational tree depths, leaf
potentials, and offsets.  We may then take, for example, a Laurent-series
field with real residue field, which is infinite.
\end{proof}

Clear the denominator of $\lambda$ at the same time as those of $p_I$; this
multiplies every valuation below by the same positive integer and does not
change any inequality.  Thus assume $p_I,\lambda\in\mathbb Z$ and normalize
$\val(\pi)=1$.

Let
\[
 L_0=\sum_e\OO v_e\subseteq\K^2.
\]
Because the support has rank two, $L_0$ is a full-rank lattice.  For a clone
subset $T$, set
\[
 L_T=L_0+\sum_{e\in T}\OO\pi^{-\lambda}v_e.
\]
The determinant ideal of $L_T$ is generated by its pairs of displayed
generators.  Hence
\[
 \val(\bigwedge^2 L_T)
 =\min_I\{p_I-\lambda|I\cap T|\}.
\]
Taking $T=\widetilde E_J$ identifies the right-hand side with
$-r_\lambda(J)$.  Smith normal form gives
\begin{equation}
 r_\lambda(J)-r_\lambda(\varnothing)
 =\length_\OO(L_{\widetilde E_J}/L_0).
 \label{eq:length-rank}
\end{equation}

All these quotients live in the single finite-length module
\[
 Q=\pi^{-\lambda}L_0/L_0.
\]
For each membership atom $\omega$, let $U_\omega\subseteq Q$ be generated by
the images of $\pi^{-\lambda}v_e$ for its two clones.  Then
\begin{equation}
 h(J):=r_\lambda(J)-r_\lambda(\varnothing)
 =\length_\OO\left(
   \sum_{\omega:\,\omega\cap J\neq\varnothing}U_\omega
  \right).
 \label{eq:module-rank}
\end{equation}

\section{Module-length Ingleton and completion of the proof}

\begin{lemma}[Module-length Ingleton]
\label{lem:module-ingleton}
Let $Q$ be a finite-length module and let $A,B,C,D$ be four submodules.  If
$h(S)$ is the length of the sum of the submodules named by $S$, then
\begin{align}
 &h(AB)+h(AC)+h(AD)+h(BC)+h(BD)\notag\\
 &\qquad\geq h(A)+h(B)+h(ABC)+h(ABD)+h(CD).
 \label{eq:ingleton}
\end{align}
\end{lemma}

\begin{proof}
Write
\[
 I(A;B)=h(A)+h(B)-h(AB)=\length(A\cap B)
\]
and define
\[
 I(A;B\mid C)=h(AC)+h(BC)-h(C)-h(ABC).
\]
Put $Z=A\cap B$.  Modularity of length gives
\begin{align*}
 I(A;B\mid C)&\geq\length(Z)-\length(Z\cap C),\\
 I(A;B\mid D)&\geq\length(Z)-\length(Z\cap D),\\
 I(C;D)&=\length(C\cap D).
\end{align*}
Moreover,
\[
 \length(Z\cap C)+\length(Z\cap D)
 \leq\length(Z)+\length(C\cap D),
\]
because the sum of the two intersections lies in $Z$ and their intersection
lies in $C\cap D$.  Thus
\[
 I(A;B\mid C)+I(A;B\mid D)+I(C;D)\geq I(A;B),
\]
which expands to \eqref{eq:ingleton}.
\end{proof}

Apply the lemma to
\[
 H_{\mathsf X}=\sum_{\omega:\,\mathsf X\in\omega}U_\omega,
 \qquad \mathsf X\in\mathcal L.
\]
These four submodules may overlap.  By \eqref{eq:module-rank},
$r_\lambda$ satisfies \eqref{eq:ingleton}, since adding the same constant to a
set function cancels from the five-against-five expression.

Combining this with \eqref{eq:evaluation-valuation},
\begin{align*}
 \val_t\mathcal R_q
 &=\sum_{J\in\{\mathsf A,\mathsf B,\mathsf{ABC},\mathsf{ABD},\mathsf{CD}\}}
     r_\lambda(J)\\
 &\quad-\sum_{J\in\{\mathsf{AB},\mathsf{AC},\mathsf{AD},
                       \mathsf{BC},\mathsf{BD}\}}r_\lambda(J)\\
 &\leq0.
\end{align*}
This contradicts \eqref{eq:positive-valuation}.  Therefore the normalized
fixed-dimensional semialgebraic family has no sequence with ratio tending to
zero.  Its infimum is positive, proving Theorem~\ref{thm:main}.

\section{The exact upper-bound family}

Use variables $a,b,c,d,o$, take $A,B,C,D$ to be the four corresponding
singletons, and leave $o$ outside their union.  For $L,x>0$, define
\begin{align}
 f_{L,x}={}&xL^{-3}(ac+bc)+\frac{x}{4}L^{-8}(ad+bd)
 +xL^{-2}(ao+bo)\notag\\
 &+L^{-1}cd+L^5co+\frac14do.
 \label{eq:upper-family}
\end{align}
It is the Cauchy--Binet expansion of
\[
 \det\left(\sum_{i\in\{a,b,c,d,o\}}z_i v_iv_i^{\mathsf T}\right)
\]
for
\begin{align*}
 v_a=v_b&=\sqrt{x}L^{-5/2}(1,0),&
 v_c&=L(0,1),\\
 v_d&=L^{-3/2}(1,1/2),&
 v_o&=L^{3/2}(1,1).
\end{align*}
The determinant of a positive-semidefinite linear pencil is real stable, and
homogeneous real stable polynomials with nonnegative coefficients are
Lorentzian \cite{BrandenHuh2020}.  Thus $f_{L,x}$ is a nonnegative
Lorentzian quadratic.

Set $\varepsilon=L^{-8}$.  The leading powers and coefficients of the ten
evaluations are
\[
\begin{array}{c|c@{\qquad}c|c}
AB&2xL^{-10}&A&xL^{-10}\\
AC&(x+1)L^{-3}&B&xL^{-10}\\
AD&(x+1)L^{-8}/4&ABC&(2x+1)L^{-3}\\
BC&(x+1)L^{-3}&ABD&(2x+1)L^{-8}/4\\
BD&(x+1)L^{-8}/4&CD&L^{-1}.
\end{array}
\]
Both products have total $L$-exponent $-32$, and therefore
\begin{equation}
 \lim_{L\to\infty}\mathcal R_{f_{L,x}}(L^{-8})
 =R_*(x)=\frac{(x+1)^4}{2x(2x+1)^2}.
 \label{eq:limit-ratio}
\end{equation}
Its logarithmic derivative has the sign of
\[
 2x^2-3x-1.
\]
Since $R_*$ diverges at both ends of the positive half-line, its unique
minimum occurs at
\[
 x_*=(3+\sqrt{17})/4.
\]
Substitution into \eqref{eq:limit-ratio} gives
\[
 R_*(x_*)=\frac{-107+51\sqrt{17}}{128}.
\]
This proves Proposition~\ref{prop:upper}.

For a finite rational check, $L=4$ and $x=16/9$ give
\[
 \mathcal R=
 \frac{1151375679760487544161993156626743807199200025995873}
 {1309813763635875854527900038158486602837796044630625}
 <\frac9{10}.
\]
The repository computes this identity using exact rational arithmetic.

\section{Separation from a common principal-minor model}

The ten subsets in \eqref{eq:ratio} have the same incidence pattern as the
classical $4\times4$ principal-minor Ingleton functional.  The following
exact example shows that Theorem~\ref{thm:main} is not obtained by simply
placing its ten values into one positive-semidefinite matrix.

For a matrix $K$ and $I\subseteq\{1,2,3,4\}$, write $K[I]$ for the principal
minor indexed by $I$.  A positive \emph{modular gauge} has the form
\[
 W(I)=\lambda\left(\prod_{i\in I}a_i\right)V(I),
 \qquad \lambda,a_1,a_2,a_3,a_4>0.
\]
The factors in such a gauge cancel from the five-against-five ratio.

\begin{proposition}[Exact model separation]
\label{prop:separation}
There are a nonnegative Lorentzian quadratic, four subsets, and a positive
regularizer for which no real symmetric positive-semidefinite $4\times4$
matrix has the ten specified principal minors equal to the ten evaluations,
even after an arbitrary positive modular gauge.
\end{proposition}

\begin{proof}
Take
\[
 f(x_0,x_1)=x_0x_1,\qquad \varepsilon=\frac13,
 \qquad A=\{0\},\quad B=\{1\},\quad C=\varnothing,
 \quad D=\{0,1\}.
\]
In the order
\[
 (AB,AC,AD,BC,BD;A,B,ABC,ABD,CD),
\]
the ten values are
\begin{equation}
 \left(
 \frac{16}{9},\frac49,\frac{16}{9},\frac49,\frac{16}{9};
 \frac49,\frac49,\frac{16}{9},\frac{16}{9},\frac{16}{9}
 \right).
 \label{eq:separation-values}
\end{equation}
The Hessian of $f$ has eigenvalues $1,-1$, so $f$ is Lorentzian.

Use the principal-minor labels $1=C,2=D,3=A,4=B$.  Put
\[
 a=V(A),\quad b=V(B),\quad u=V(AB),\quad
 P=V(ABC),\quad Q=V(ABD)
\]
and form the five gauge invariants
\begin{align*}
 x&=\frac{uV(AC)}{aP},&
 y&=\frac{uV(BC)}{bP},&
 z&=\frac{uV(AD)}{aQ},&
 w&=\frac{uV(BD)}{bQ},&
 h&=\frac{u^3V(CD)}{abPQ}.
\end{align*}
The five positive gauge parameters normalize
$V(A),V(B),V(AB),V(ABC),V(ABD)$ to $1$; the remaining five coordinates are
$(x,y,z,w,h)$.  Existence of a representation is invariant under this
normalization.  Any remaining four diagonal gauge factors can then be
absorbed by a positive diagonal congruence of the representing matrix, while
the common scalar remains.  Thus a modular-gauge principal-minor
representation of \eqref{eq:separation-values} would yield a
positive-semidefinite matrix $K$ and a scalar $\lambda>0$ satisfying
\begin{gather}
 K[3]=K[4]=K[13]=K[14]=K[34]=K[134]=K[234]=\lambda,
 \label{eq:canonical-minors-one}\\
 K[23]=K[24]=4\lambda,\qquad K[12]=16\lambda,
 \label{eq:canonical-minors-two}
\end{gather}
because the exact invariants are $(1,1,4,4,16)$.

Let $R=K_{\{3,4\},\{3,4\}}$.  Equations
\eqref{eq:canonical-minors-one} give
\[
 R=\begin{pmatrix}\lambda&r\\r&\lambda\end{pmatrix},
 \qquad \det R=\lambda,
 \qquad r^2=\lambda^2-\lambda.
\]
Hence $\lambda\geq1$ and $R$ is positive definite.  For $i=1,2$, write
\[
 K_{i,\{3,4\}}=R\beta_i,
 \qquad s_i=K_{ii}-\beta_i^{\mathsf T}R\beta_i.
\]
Schur complements give $s_1=s_2=1$.  If
$\beta_i=((\beta_i)_3,(\beta_i)_4)$, direct expansion of the two-by-two
minors gives
\[
 K[i3]=\lambda s_i+\lambda(\beta_i)_4^2,
 \qquad
 K[i4]=\lambda s_i+\lambda(\beta_i)_3^2.
\]
It follows that $\beta_1=0$ and $K_{11}=1$, while
$| (\beta_2)_3|^2=|(\beta_2)_4|^2=3$.  Therefore
\[
 K_{22}=1+\beta_2^{\mathsf T}R\beta_2
 \leq1+6\lambda+6|r|
 \leq1+12\lambda.
\]
For every value of the unconstrained entry $K_{12}$,
\[
 K[12]=K_{11}K_{22}-K_{12}^2
 \leq1+12\lambda<16\lambda,
\]
contradicting \eqref{eq:canonical-minors-two}.  Notice that the other five
principal minors were left completely free.
\end{proof}

Proposition~\ref{prop:separation} rules out the natural simultaneous
common-matrix reduction, not every possible indirect use of a
principal-minor inequality.  Its exact arithmetic is included in the
certificate manifest.

\section{Relation to adjacent results and open problems}

The closest classical input is additive Ingleton for linear rank functions
\cite{Ingleton1971}.  The proof above needs a valuated form because ordinary
support ranks do not record relative coefficient-vanishing rates.

Hall and Johnson studied bounded ratios of products of principal minors
\cite{HallJohnson2008}.  Boege and Bouthat recently proved that the infimum of
the Ingleton ratio over $4\times4$ positive definite matrices is $16/27$;
Sendov independently proved the reciprocal sharp supremum formulation
\cite{BoegeBouthat2026,Sendov2026}.  Their domain and evaluations are different from
those of Theorem~\ref{thm:main}.  After relabelling, the ten subsets follow
the same combinatorial Ingleton pattern, so this is the main collision risk.
Here, however, the input is an arbitrary nonnegative Lorentzian quadratic in
any number of variables, the four sets may overlap, and the ten quantities are
regularized polynomial evaluations.  No simultaneous representation of
these ten values as scalar principal minors of one fixed $4\times4$ matrix
exists in general, even up to a positive ratio-neutral modular rescaling, by
Proposition~\ref{prop:separation}.  The proposition does not exclude a more
indirect implication.  Neither result is asserted to imply the other.

Al Ahmadieh, Rinc\'on, Vinzant, and Yu characterized tropicalized principal
minors and developed coefficient inequalities valid for Lorentzian
polynomials \cite{AlAhmadiehRinconVinzantYu2025}.  In our dated statement-level
search, no theorem there matched the regularized ten-evaluation bound above.

The following questions remain open.
\begin{enumerate}
\item Determine the optimal value of $c_2^\star$, or any explicit positive lower
bound.
\item Decide whether an analogous degree-only bound holds for Lorentzian
polynomials of degree three.
\item Identify natural subclasses for which the sharp constant is $1$.
\end{enumerate}

The repository's dated novelty audit records the databases and queries used
to search for a matching theorem.  Failure to find a match is evidence about
that search, not a mathematical proof of priority.

\paragraph{Verification provenance.}
The argument was developed with computer-assisted exploration and repeated
adversarial proof review.  The exact-arithmetic programs in the repository
check conventions, the finite witness, the principal-minor separation, and a
bounded rank-two valuation grid; no finite enumeration is used in place of the
general proof above.

\bibliographystyle{alpha}
\bibliography{references}

\end{document}
